% Vietnamese translation for OI Public Library
% Source-PDF: IOI中国国家候选队论文/2001/符文杰.pdf
\documentclass[11pt]{article}
\usepackage{oipl-vn}

\title{Nguyên Lý Pólya Và Ứng Dụng}
\author{Fu Wenjie}
\date{2001}

\begin{document}
\maketitle

\origtitle{Pólya principle and applications}
\sourcepath{IOI Candidate Team Papers/2001/Fu Wenjie.pdf}

\begin{abstract}
Nguyên lý Pólya là một công cụ hiệu quả và gọn trong tổ hợp, dùng để đếm số
trạng thái tổ hợp khác nhau về bản chất khi có các phép đối xứng. Bài viết
giới thiệu các khái niệm nhóm, hoán vị, nhóm hoán vị, bổ đề Burnside, định lý
Pólya, rồi áp dụng vào bài toán tô màu hình vuông dưới phép quay.
\end{abstract}

\section{Ví dụ mở đầu}

\paragraph{Ví dụ 1.}
Tô màu một bảng \(2\times2\) bằng hai màu đen và trắng. Hỏi có thể thu được
bao nhiêu hình khác nhau? Hai phương án được xem là cùng một phương án nếu có
thể quay để trùng nhau.

Vì quy mô bài toán rất nhỏ, ta có thể liệt kê toàn bộ các phương án tô màu.
Trong bản gốc, mười sáu cách tô được ký hiệu là \(C_1,\ldots,C_{16}\). Có
bốn phép quay của bảng \(2\times2\): quay \(0^\circ\), \(90^\circ\),
\(180^\circ\), và \(270^\circ\). Trong bài viết này, phép quay mặc định là
quay theo chiều kim đồng hồ.

Sau khi thử, ta thấy chỉ có \(6\) phương án khác nhau về bản chất:
\[
  \{C_1\},\quad \{C_2\},\quad
  \{C_3,C_4,C_5,C_6\},\quad
  \{C_7,C_8,C_9,C_{10}\},\quad
  \{C_{11},C_{12}\},\quad
  \{C_{13},C_{14},C_{15},C_{16}\}.
\]
Nói cách khác, các phần tử trong cùng một tập trên có thể biến đổi qua lại
bằng phép quay và phải được tính là một phương án.

Để minh họa cách đánh số, ta dùng \(0\) cho trắng và \(1\) cho đen, đọc bốn ô
theo thứ tự trái trên, phải trên, trái dưới, phải dưới. Khi đó mười sáu cách
tô có thể xem là các xâu nhị phân từ \(0000\) tới \(1111\). Hình trong bản
gốc liệt kê các hình này và sáu lớp tương đương; ở đây ta dùng mô tả bằng ký
hiệu vì hình không còn khôi phục được từ văn bản trích xuất.

Khi bài toán đổi từ bảng \(2\times2\) thành bảng \(20\times20\), thậm chí
\(200\times200\), ta không thể liệt kê từng phương án nữa. Khi đó nguyên lý
Pólya trở thành một cách giải rất tốt. Trước khi đi vào nguyên lý Pólya, ta
cần nhắc lại một vài khái niệm.

\section{Nhóm, hoán vị và nhóm hoán vị}

\paragraph{Nhóm.}
Cho một tập \(G=\{a,b,c,\ldots\}\) và một phép toán hai ngôi trên \(G\). Nếu
thỏa mãn bốn điều kiện sau:
\begin{enumerate}
  \item \term{Đóng}: với mọi \(a,b\in G\), tồn tại \(c\in G\) sao cho
  \(a*b=c\).
  \item \term{Kết hợp}: với mọi \(a,b,c\in G\),
  \[
    (a*b)*c=a*(b*c).
  \]
  \item \term{Phần tử đơn vị}: tồn tại \(e\in G\) sao cho với mọi \(a\in G\),
  \[
    a*e=e*a=a.
  \]
  \item \term{Phần tử nghịch đảo}: với mọi \(a\in G\), tồn tại \(b\in G\) sao
  cho
  \[
    a*b=b*a=e.
  \]
  Ta ký hiệu \(b=a^{-1}\).
\end{enumerate}
thì \(G\) được gọi là một nhóm dưới phép toán \(*\). Thường viết gọn \(a*b\)
thành \(ab\).

\paragraph{Hoán vị.}
Một hoán vị trên \(n\) phần tử \(1,2,\ldots,n\) được viết
\[
  \begin{pmatrix}
    1 & 2 & \cdots & n\\
    a_1 & a_2 & \cdots & a_n
  \end{pmatrix},
\]
nghĩa là \(1\) được thay bởi một số \(a_1\) trong \(1..n\), \(2\) được thay
bởi một số \(a_2\), và cứ thế tới \(n\), trong đó
\(a_1,a_2,\ldots,a_n\) đôi một khác nhau.

Trong ví dụ tô màu \(2\times2\), xét mười sáu phương án \(C_1,\ldots,C_{16}\)
như các phần tử đang bị tác động bởi phép quay. Bốn phép quay tạo ra bốn hoán
vị sau:
\[
\begin{array}{c|rrrrrrrrrrrrrrrr}
 &1&2&3&4&5&6&7&8&9&10&11&12&13&14&15&16\\ \hline
a_1&1&2&3&4&5&6&7&8&9&10&11&12&13&14&15&16\\
a_2&1&2&6&3&4&5&10&7&8&9&12&11&16&13&14&15\\
a_3&1&2&5&6&3&4&9&10&7&8&11&12&15&16&13&14\\
a_4&1&2&4&5&6&3&8&9&10&7&12&11&14&15&16&13
\end{array}
\]
Ở đây \(a_1,a_2,a_3,a_4\) lần lượt là quay \(0^\circ,90^\circ,180^\circ\)
và \(270^\circ\).

\paragraph{Nhóm hoán vị.}
Nhóm hoán vị là một nhóm mà phần tử là các hoán vị, còn phép toán là phép nối
hoán vị. Ví dụ:
\[
\begin{pmatrix}
1&2&3&4\\
3&1&2&4
\end{pmatrix}
\begin{pmatrix}
1&2&3&4\\
4&3&2&1
\end{pmatrix}
=
\begin{pmatrix}
1&2&3&4\\
2&4&3&1
\end{pmatrix}.
\]
Có thể kiểm tra rằng nhóm hoán vị thỏa mãn bốn điều kiện của nhóm. Trong ví
dụ đang xét,
\[
  G=\{\text{quay }0^\circ,\text{quay }90^\circ,
       \text{quay }180^\circ,\text{quay }270^\circ\}.
\]

\section{Bổ đề Burnside}

Ta xét công thức
\[
  |E_k|\cdot |Z_k|=|G|,\qquad k=1,\ldots,n.
\]
Công thức này liên quan trực tiếp tới số phần tử của nhóm và rất hữu ích.

\paragraph{Lớp hoán vị giữ cố định \(K\).}
Giả sử \(G\) là một nhóm hoán vị trên \(1..n\). Nếu \(K\) là một phần tử trong
\(1..n\), tập tất cả các hoán vị trong \(G\) làm \(K\) không đổi được ký hiệu
là \(Z_K\), gọi là lớp hoán vị trong \(G\) giữ cố định \(K\).

Trong ví dụ này, \(G\) là nhóm hoán vị trên các phương án tô màu
\(1,\ldots,16\). Với phương án \(1\), cả bốn hoán vị đều giữ nó không đổi, nên
\[
  Z_1=\{a_1,a_2,a_3,a_4\}.
\]
Với phương án \(3\), chỉ \(a_1\) giữ nó không đổi, nên \(Z_3=\{a_1\}\). Với
phương án \(11\), hai hoán vị \(a_1\) và \(a_3\) giữ nó không đổi, nên
\[
  Z_{11}=\{a_1,a_3\}.
\]

\paragraph{Lớp tương đương.}
Giả sử \(G\) là một nhóm hoán vị trên \(1..n\). Nếu \(K\) là một phần tử trong
\(1..n\), quỹ đạo của \(K\) dưới tác động của \(G\) được ký hiệu là \(E_K\).
Đó là tập tất cả phần tử mà \(K\) có thể biến thành khi áp dụng các hoán vị
trong \(G\).

Trong ví dụ:
\[
  E_1=\{1\},\qquad E_3=\{3,4,5,6\},\qquad E_{11}=\{11,12\}.
\]
Các dữ liệu này khớp hoàn toàn với công thức:
\[
  |E_1||Z_1|=1\cdot4=4=|G|,
\]
\[
  |E_3||Z_3|=4\cdot1=4=|G|,
\]
\[
  |E_{11}||Z_{11}|=2\cdot2=4=|G|.
\]
Do giới hạn dung lượng, bản gốc không chứng minh định lý này.

Bây giờ xét tổng số lần mỗi phần tử được giữ cố định bởi các hoán vị. Đặt
\[
  S_{ij}=
  \begin{cases}
    0, & a_i\notin Z_j,\ \text{nghĩa là } j \text{ bị thay đổi dưới } a_i,\\
    1, & a_i\in Z_j,\ \text{nghĩa là } j \text{ không đổi dưới } a_i.
  \end{cases}
\]
Ký hiệu \(D(a_i)\) là số phần tử không đổi dưới hoán vị \(a_i\). Bảng tổng
quát có dạng:
\[
\begin{array}{c|ccccc|c}
\text{hoán vị} & 1 & 2 & 3 & \cdots & 16 & D(a_i)\\ \hline
a_1 & S_{1,1} & S_{1,2} & S_{1,3} & \cdots & S_{1,16} & D(a_1)\\
a_2 & S_{2,1} & S_{2,2} & S_{2,3} & \cdots & S_{2,16} & D(a_2)\\
a_3 & S_{3,1} & S_{3,2} & S_{3,3} & \cdots & S_{3,16} & D(a_3)\\
a_4 & S_{4,1} & S_{4,2} & S_{4,3} & \cdots & S_{4,16} & D(a_4)\\ \hline
|Z_j| & |Z_1| & |Z_2| & |Z_3| & \cdots & |Z_{16}| &
\end{array}
\]
Ví dụ: phương án \(1\) không đổi dưới \(a_1\), nên \(S_{1,1}=1\). Phương án
\(3\) bị thay đổi dưới \(a_3\), nên \(S_{3,3}=0\). Dưới \(a_1\), cả \(16\)
phương án đều không đổi, nên \(D(a_1)=16\). Dưới \(a_2\), chỉ có hai phương
án \(1\) và \(2\) không đổi, nên \(D(a_2)=2\).

Trong trường hợp tổng quát, ta có
\[
  \sum_{j=1}^{n}|Z_j|=\sum_{i=1}^{s}D(a_i),
\]
trong đó \(s=|G|\). Xét vế trái. Giả sử
\[
  N=\{1,\ldots,n\}
\]
có \(L\) lớp tương đương:
\[
  N=E_1+E_2+\cdots+E_L.
\]
Nếu \(j\) và \(k\) thuộc cùng một lớp tương đương thì \(|Z_j|=|Z_k|\). Do đó
\[
  \sum_{k=1}^{n}|Z_k|
  =\sum_{i=1}^{L}\sum_{k\in E_i}|Z_k|
  =\sum_{i=1}^{L}|E_i|\cdot |Z_i|
  =L\cdot |G|.
\]
Ở đây \(L\) chính là số trạng thái tổ hợp khác nhau về bản chất mà ta cần
tìm. Vì vậy
\[
  L=\frac{1}{|G|}\sum_{k=1}^{n}|Z_k|
   =\frac{1}{|G|}\sum_{j=1}^{s}D(a_j).
\]
Áp dụng vào ví dụ đầu, ta được
\[
  L=\frac{16+2+4+2}{4}=6,
\]
đúng với kết quả liệt kê. Công thức này gọi là \term{bổ đề Burnside}.

Tuy nhiên, việc tính \(D(a_j)\) không hề dễ nếu làm trực tiếp. Nếu dùng tìm
kiếm, độ phức tạp tổng quát là
\[
  \mathcal{O}(nsp),
\]
trong đó \(n\) là số phần tử, \(s\) là số hoán vị, còn \(p\) là số ô. Ở đây
\(n\) có thể rất lớn. Vì vậy bước tiếp theo là tìm một cách tính \(D(a_j)\)
đơn giản hơn.

\section{Chu trình và định lý Pólya}

\paragraph{Chu trình.}
Ký hiệu
\[
  (a_1a_2\cdots a_n)=
  \begin{pmatrix}
    a_1&a_2&\cdots&a_{n-1}&a_n\\
    a_2&a_3&\cdots&a_n&a_1
  \end{pmatrix}.
\]
Đây gọi là một chu trình bậc \(n\). Mọi hoán vị đều có thể viết thành tích
của một số chu trình đôi một rời nhau. Hai chu trình
\((a_1a_2\cdots a_n)\) và \((b_1b_2\cdots b_m)\) được gọi là rời nhau nếu
\[
  a_i\ne b_j
\]
với mọi \(i,j\). Ví dụ:
\[
  \begin{pmatrix}
    1&2&3&4&5\\
    3&5&1&4&2
  \end{pmatrix}
  =(13)(25)(4).
\]
Cách biểu diễn này là duy nhất nếu bỏ qua thứ tự viết các chu trình. Số chu
trình của hoán vị là số chu trình xuất hiện trong biểu diễn đó. Chẳng hạn
\((13)(25)(4)\) có \(3\) chu trình.

Với các khái niệm trên, ta đổi góc nhìn cho bài toán \(2\times2\). Đánh số
bốn ô của bảng như sau:
\[
\begin{array}{|c|c|}
\hline
2&1\\ \hline
3&4\\ \hline
\end{array}
\]
Xây dựng nhóm hoán vị
\[
  G'=\{g_1,g_2,g_3,g_4\},\qquad |G'|=4,
\]
và gọi \(c(g_i)\) là số chu trình của \(g_i\). Dưới tác động của \(G'\), ta
có:
\[
\begin{array}{c|c|c}
\text{phép quay} & \text{dạng chu trình} & c(g_i)\\ \hline
g_1:\ 0^\circ & (1)(2)(3)(4) & 4\\
g_2:\ 90^\circ & (4\ 3\ 2\ 1) & 1\\
g_3:\ 180^\circ & (1\ 3)(2\ 4) & 2\\
g_4:\ 270^\circ & (1\ 2\ 3\ 4) & 1
\end{array}
\]

Ta nhận thấy: số hình nhận được khi các đối tượng trong cùng một chu trình
của \(g_i\) phải tô cùng màu là \(m^{c(g_i)}\), và con số này đúng bằng số
hình không đổi dưới hoán vị \(a_i\) trong nhóm \(G\). Với \(m=2\):
\[
  2^{c(g_1)}=2^4=16=D(a_1),
\]
\[
  2^{c(g_2)}=2^1=2=D(a_2),
\]
\[
  2^{c(g_3)}=2^2=4=D(a_3),
\]
\[
  2^{c(g_4)}=2^1=2=D(a_4).
\]

Từ đó rút ra kết luận sau.

\paragraph{Định lý Pólya.}
Giả sử \(G=\{g_1,\ldots,g_s\}\) là một nhóm hoán vị trên \(p\) đối tượng. Tô
\(p\) đối tượng bằng \(m\) màu. Khi đó số phương án tô màu khác nhau về bản
chất là
\[
  L=\frac{1}{|G|}
    \left(m^{c(g_1)}+m^{c(g_2)}+\cdots+m^{c(g_s)}\right),
\]
trong đó \(c(g_i)\) là số chu trình của hoán vị \(g_i\).

Đây chính là định lý Pólya. Nếu dùng định lý Pólya, độ phức tạp là
\[
  \mathcal{O}(sp),
\]
với \(s\) là số hoán vị và \(p\) là số ô. So với cách dùng bổ đề Burnside một
cách trực tiếp, đây là một cải thiện lớn. Định lý Pólya khai thác cấu trúc
nội tại của đối tượng, tổng kết quy luật, loại bỏ nhiều tìm kiếm mù quáng
không cần thiết, và hạ độ phức tạp của lớp bài toán này xuống mức rất thấp.

\section{Bài toán bảng vuông tổng quát}

Bây giờ đổi bài toán thành: với bảng \(n\times n\), mỗi ô có thể tô một trong
\(m\) màu, hãy tính số phương án khác nhau về bản chất dưới các phép quay.

\subsection{Phân tích bằng tìm kiếm cải tiến}

Trước hết xét cách tìm kiếm dễ nghĩ tới nhất, được đưa trong phụ lục. Cách
tìm kiếm đó có hiệu quả rất thấp và còn nhiều chỗ có thể cải tiến. Trước đó
ta đã dùng cách ``tìm rồi kiểm tra'', nên tính mù quáng rất cao. Điều mong
muốn hơn là vừa tìm vừa kiểm tra, tránh quá nhiều liệt kê không cần thiết,
thậm chí đưa điều kiện kiểm tra vào ngay biên tìm kiếm để loại bỏ hoàn toàn
các liệt kê vô hiệu. Mong muốn này có thể thực hiện được.

Ta có thể chia bảng thành bốn hình chữ nhật không chồng nhau, mỗi hình có
chiều dài
\[
  \left\lfloor\frac{n+1}{2}\right\rfloor
\]
và chiều rộng
\[
  \left\lfloor\frac{n}{2}\right\rfloor.
\]
Khi \(n\) chẵn, bốn hình chữ nhật này phủ kín bảng. Khi \(n\) lẻ, còn lại ô
chính giữa; ô này không đổi dưới mọi phép quay, nên tô màu nào cũng không ảnh
hưởng tới các ô khác.

Gọi \(m\) màu là \(0,1,\ldots,m-1\). Trong mỗi hình chữ nhật, ta đọc màu của
các ô theo một thứ tự cố định: hình chữ nhật trái trên đọc từ trái sang phải,
từ trên xuống dưới; hình chữ nhật phải trên đọc từ trên xuống dưới, từ phải
sang trái; và tương tự cho hai phần còn lại. Như vậy mỗi hình chữ nhật sinh
ra một số ở hệ cơ số \(m\), rồi có thể xem như một số thập phân trong khoảng
\[
  0\ \text{tới}\ m^{\lfloor n^2/4\rfloor}-1,
\]
vì
\[
  \left\lfloor\frac n2\right\rfloor
  \left\lfloor\frac{n+1}{2}\right\rfloor
  =\left\lfloor\frac{n^2}{4}\right\rfloor.
\]
Nhờ đó, một bảng được giản lược thành bốn số nguyên. Trong mỗi lớp tương
đương, chỉ cần giữ phương án có số ở góc trái trên lớn nhất; nếu bằng nhau
thì so tiếp theo chiều kim đồng hồ. Khi liệt kê, ba số còn lại không được lớn
hơn số ở góc trái trên, nên hiệu quả là tốt nhất theo cách nhìn này.

Xét kỹ hơn, khi số ở góc trái trên là \(i\), với \(0\le i\le R-1\) và
\[
  R=m^{\lfloor n^2/4\rfloor},
\]
ta có bốn loại:
\begin{enumerate}
  \item Ba số còn lại đều nhỏ hơn \(i\): có \(i^3\) phương án.
  \item Số ở góc phải trên bằng \(i\), hai số còn lại nhỏ hơn \(i\): có
  \(i^2\) phương án.
  \item Số ở góc phải trên và góc phải dưới đều bằng \(i\), số ở góc trái
  dưới không lớn hơn \(i\): có \(i+1\) phương án.
  \item Số ở góc phải dưới bằng \(i\), hai số còn lại nhỏ hơn \(i\), và số ở
  góc phải trên không nhỏ hơn số ở góc trái dưới: có \(i(i+1)/2\) phương án.
\end{enumerate}
Do đó
\[
\begin{aligned}
L
&=\sum_{i=0}^{R-1}\left(i^3+i^2+i+1+\frac12 i(i+1)\right)\\
&=\sum_{i=0}^{R-1}\left(i^3+\frac32 i^2+\frac32 i+1\right)\\
&=\sum_{i=1}^{R}\left((i-1)^3+\frac32(i-1)^2+\frac32(i-1)+1\right)\\
&=\sum_{i=1}^{R}\left(i^3-\frac32 i^2+\frac32 i\right)\\
&=\frac14 R^2(R+1)^2-\frac32\cdot\frac16R(R+1)(2R+1)
  +\frac32\cdot\frac12R(R+1)\\
&=\frac14\left(R^4+R^2+2R\right).
\end{aligned}
\]
Khi \(n\) lẻ, còn cần nhân thêm một thừa số \(m\) cho ô trung tâm.

Như vậy ta đã khéo léo thu được một công thức. Nhưng để nghĩ ra công thức
này cần nhiều trí tuệ và thời gian. Mặt khác, phương pháp này chỉ hữu dụng
cho đúng bài toán này, khó áp dụng rộng cho cả một lớp bài. Vì vậy nếu nắm
được một nguyên lý có phạm vi dùng rộng hơn, việc giải các bài tương tự sẽ
dễ hơn nhiều.

\subsection{Giải bằng định lý Pólya}

Ta dùng định lý Pólya theo ba bước.

\paragraph{Bước 1: xác định nhóm hoán vị.}
Ở đây rõ ràng chỉ có bốn hoán vị: quay \(0^\circ\), quay \(90^\circ\), quay
\(180^\circ\), quay \(270^\circ\). Vì vậy
\[
  G=\{\text{quay }0^\circ,\text{quay }90^\circ,
       \text{quay }180^\circ,\text{quay }270^\circ\}.
\]

\paragraph{Bước 2: tính số chu trình.}
Trước hết đánh số tuần tự các ô từ \(1\) tới \(n^2\), sau đó dùng một mảng
hai chiều ghi trạng thái sau phép hoán vị. Cuối cùng, tìm kiếm trên các ô để
tính số chu trình của từng phép quay. Độ phức tạp là tuyến tính theo số ô.

\paragraph{Bước 3: thay vào công thức.}
Dùng trực tiếp định lý Pólya:
\[
  L=\frac{1}{|G|}
    \left(m^{c(g_1)}+m^{c(g_2)}+\cdots+m^{c(g_s)}\right).
\]

\paragraph{Chương trình minh họa.}
Bản gốc dùng Pascal. Dưới đây giữ nguyên cấu trúc chương trình, dịch chú
thích sang tiếng Việt không dấu.

\begin{verbatim}
const
 maxn=10;
var
 a,b:array[1..maxn,1..maxn] of integer; {ghi trang thai bang}
 i,j,m,n:integer; {m la so mau; n la kich thuoc bang}
 l,l1:longint;

procedure xz; {quay bang 90 do}
var
 i,j:integer;
begin
 for i:=1 to n do
   for j:=1 to n do
     a[j,n+1-i]:=b[i,j];
 b:=a
end;

procedure xhj; {tinh so chu trinh cua trang thai hien tai}
var
 i,j,i1,j1,k,p:integer;
begin
 k:=0; {dem so chu trinh}
 for i:=1 to n do
  for j:=1 to n do
    if a[i,j]>0 then {tim o chua tham}
     begin
       inc(k); {them mot chu trinh}
       i1:=(a[i,j]-1) div n;
       j1:=(a[i,j]-1) mod n+1; {o ke tiep trong chu trinh}

       a[i,j]:=0; {danh dau da tham}
       while a[i1,j1]>0 do begin
        p:=a[i1,j1]; {luu thong tin o hien tai}
        a[i1,j1]:=0; {danh dau da tham}
        i1:=(p-1) div n+1;
        j1:=(p-1) mod n+1 {o ke tiep trong chu trinh}
       end {ket thuc khi da tham tron chu trinh}
     end;
 l1:=1;
 for i:=1 to k do l1:=l1*m; {tinh m mu k}
 l:=l+l1 {cong vao tong}
end;

begin
 writeln('Input m,n=');
 readln(m,n); {nhap du lieu}
 for i:=1 to n do
  for j:=1 to n do
    a[i,j]:=(i-1)*n+j; {khoi tao trang thai bang}
 b:=a;
 xhj; {quay 0 do}
 xz; xhj; {quay 90 do}
 xz; xhj; {quay 180 do}
 xz; xhj; {quay 270 do}
 l:=l div 4;
 writeln(l); {xuat ket qua}
 readln
end.
\end{verbatim}

Trong chương trình trên, bản gốc tạm thời tránh xử lý số lớn vì điều đó không
liên quan trực tiếp tới nội dung đang trình bày.

Nếu suy nghĩ thêm, ta còn có thể tối ưu lời giải này bằng cách xét riêng các
trường hợp của \(n\):
\begin{itemize}
  \item Khi \(n\) chẵn: dưới quay \(0^\circ\), có \(n^2\) chu trình; dưới
  quay \(180^\circ\), có \(n^2/2\) chu trình; dưới quay \(90^\circ\) và
  \(270^\circ\), có \(n^2/4\) chu trình.
  \item Khi \(n\) lẻ: dưới quay \(0^\circ\), có \(n^2\) chu trình; dưới quay
  \(180^\circ\), có \((n^2+1)/2\) chu trình; dưới quay \(90^\circ\) và
  \(270^\circ\), có \((n^2+3)/4\) chu trình.
\end{itemize}
Gộp lại:
\[
  c(0^\circ)=n^2,\qquad
  c(180^\circ)=\left\lfloor\frac{n^2+1}{2}\right\rfloor,
\]
\[
  c(90^\circ)=c(270^\circ)
  =\left\lfloor\frac{n^2+3}{4}\right\rfloor.
\]
Vì vậy
\[
  L=\frac14\left(
    m^{n^2}
    +m^{\left\lfloor (n^2+3)/4\right\rfloor}
    +m^{\left\lfloor (n^2+1)/2\right\rfloor}
    +m^{\left\lfloor (n^2+3)/4\right\rfloor}
  \right).
\]
Công thức này khớp hoàn toàn với kết quả thu được bằng phân tích tìm kiếm ở
trên. Sau phần phân tích, một bài tưởng như khó đã được giải gọn; phần còn
lại trong lập trình chỉ là xử lý số lớn.

Qua các ví dụ trên, có thể thấy ưu thế của nguyên lý Pólya rất rõ khi so với
tìm kiếm. Nó không chỉ nâng cao mạnh hiệu quả thời gian của chương trình, mà
thậm chí còn không làm tăng độ khó lập trình. Trí tuệ và kinh nghiệm rất quan
trọng, nhưng nắm được các nguyên lý nền tảng giúp việc giải bài chắc chắn hơn.
Vì vậy, bên cạnh giải bài, ta cũng không nên quên học các kiến thức có tính
nguyên lý.

\section{Phụ lục: chương trình tìm kiếm}

Chương trình sau là phương pháp tìm kiếm được nhắc tới ở phần phân tích. Nó
liệt kê mọi cách tô và chỉ đếm phương án lớn nhất trong mỗi lớp tương đương
dưới phép quay.

\begin{verbatim}
const
 maxn=10;
type
 sqtype=array[1..maxn,1..maxn] of byte;
var
 n,total,m:integer;
 sq:sqtype;

function big:boolean; {kiem tra phuong an co lon nhat trong lop khong}
var
 units:array[2..4] of sqtype; {ba trang thai sau khi quay}
 i,j,k:integer;
begin
 for i:=1 to n do
   for j:=1 to n do begin
     units[2,j,n+1-i]:=sq[i,j];
     units[3,n+1-i,n+1-j]:=sq[i,j];
     units[4,n+1-j,i]:=sq[i,j]
   end;
 big:=false;
 for k:=2 to 4 do {so sanh}
  for i:=1 to n do begin
   j:=1;
   while (j<=n)and(sq[i,j]=units[k,i,j]) do inc(j);
   if j<=n then
     if sq[i,j]<units[k,i,j]
       then exit
       else break
  end;
 big:=true
end;

procedure make(x,y:byte); {liet ke mau cho tung o}
var i:integer;
begin
 if x>n then begin
   if big then inc(total);
   exit
 end;
 for i:=1 to m do begin
   sq[x,y]:=i;
   if y=n then make(x+1,1) else make(x,y+1)
 end
end;

begin
 writeln('Input m,n=');
 readln(m,n);
 total:=0;
 make(1,1);
 writeln(total);
 readln
end.
\end{verbatim}

\end{document}
